\begin{thebibliography}{99}
\addcontentsline{toc}{chapter}{References}
\singlespacing
\bibitem{ref1} 3GPP, "Evolved Universal Terrestrial Radio Access (E-UTRA); Radio Resource Control (RRC); Protocol specification," 3GPP TS 36.331, Release 17, 2022.

\bibitem{ref2} 3GPP, "Evolved Universal Terrestrial Radio Access (E-UTRA) and Evolved Universal Terrestrial Radio Access Network (E-UTRAN); Overall description; Stage 2," 3GPP TS 36.300, Release 17, 2022.

\bibitem{ref3} H. Deng, C. Peng, A. Fida, J. Meng, and Y. C. Hu, "Mobility support in cellular networks: a measurement study on its configurations and implications," in Proc. ACM Internet Measurement Conference (IMC), 2018, doi:10.1145/3278532.3278546.

\bibitem{ref4} M. Ghoshal, I. Khan, P. Dinh, Z. J. Kong, O. Basit, S. Wang, Y. Feng, Y. C. Hu, and D. Koutsonikolas, "Handover configurations in operational 5G networks: diversity, evolution, and impact on performance," arXiv:2511.03116, 2025.

\bibitem{ref5} A. Hassan, A. Narayanan, A. Zhang, W. Ye, R. Zhu, S. Jin, J. Carpenter, Z. M. Mao, F. Qian, and Z.-L. Zhang, "Vivisecting mobility management in 5G cellular networks," in Proc. ACM SIGCOMM, 2022, pp. 86--100, doi:10.1145/3544216.3544217.

\bibitem{ref6} Z. Liu, Q. Deng, Z. Tan, Z. Qian, X. Zhang, A. Swami, and S. V. Krishnamurthy, "M2HO: Mitigating the adverse effects of 5G handovers on TCP," in Proc. ACM MobiCom, 2024, doi:10.1145/3636534.3690680.

\bibitem{ref7} T. Ankome and E. Hanada, "Reactive to predictive mobility management: a systematic review of ML-driven handover optimization in 5G and beyond," Machine Learning and Knowledge Extraction, vol. 8, no. 5, art. 133, 2026, doi:10.3390/make8050133.

\bibitem{ref8} B. Saoud, I. Shayea, M. A. Alnakhli, and H. Mohamad, "Mobility and handover management in 5G/6G networks: challenges, innovations, and sustainable solutions," Technologies, vol. 13, no. 8, art. 352, 2025, doi:10.3390/technologies13080352.

\bibitem{ref9} C. Chabira, I. Shayea, G. Nurzhaubayeva, L. Aldasheva, D. Yedilkhan, and S. Amanzholova, "AI-driven handover management and load balancing optimization in ultra-dense 5G/6G cellular networks," Technologies, vol. 13, no. 7, art. 276, 2025, doi:10.3390/technologies13070276.

\bibitem{ref10} K. Nguyen Dang Dinh, P. Fazio, and M. Voznak, "A proposed mobility and resource prediction method for seamless handover and service continuity in 5G small cell networks," PLOS ONE, vol. 21, no. 8, art. e0355372, 2026, doi:10.1371/journal.pone.0355372.

\bibitem{ref11} J. M. S\'anchez-Mart\'{\i}n, J. L. Bejarano-Luque, C. Gij\'on, M. Toril, and S. Luna-Ram\'{\i}rez, "GRIMCELL: a graph neural network to predict the impact of new cells in mobile networks," Machine Learning and Knowledge Extraction, vol. 8, no. 9, art. 260, 2026, doi:10.3390/make8090260.

\bibitem{ref12} N. Mehregan and R. E. De Grande, "GCN-based throughput-oriented handover management in dense 5G vehicular networks," arXiv:2505.04894, IEEE DCOSS-IoT, 2025.

\bibitem{ref13} A. Graser, A. Jalali, J. Lampert, A. Wei\ss{}enfeld, and K. Janowicz, "MobilityDL: a review of deep learning from trajectory data," GeoInformatica, vol. 29, no. 1, pp. 115--147, 2025, doi:10.1007/s10707-024-00518-8.

\bibitem{ref14} D. Zidic, T. Mastelic, I. Nizetic Kosovic, M. Cagalj, and J. Lorincz, "Analyses of ping-pong handovers in real 4G telecommunication networks," Computer Networks, vol. 227, art. 109699, 2023.

\bibitem{ref15} A. Amirova, A. Abdiraman, L. Aldasheva, I. Shayea, D. Yedilkhan, and A. Tussupov, "Data-driven and machine learning-based analysis of handover behavior and network stability in mobile networks," Future Internet, vol. 18, no. 6, art. 290, 2026, doi:10.3390/fi18060290.

\bibitem{ref16} S. Wiegrebe, P. Kopper, R. Sonabend, B. Bischl, and A. Bender, "Deep learning for survival analysis: a review," Artificial Intelligence Review, vol. 57, no. 3, art. 65, 2024, doi:10.1007/s10462-023-10681-3.

\bibitem{ref17} H.-C. Thorsen-Meyer, D. Placido, B. S. Kaas-Hansen, et al., "Discrete-time survival analysis in the critically ill: a deep learning approach using heterogeneous data," npj Digital Medicine, vol. 5, art. 142, 2022, doi:10.1038/s41746-022-00679-6.

\bibitem{ref18} E. J. Cand\`es, L. Lei, and Z. Ren, "Conformalized survival analysis," Journal of the Royal Statistical Society Series B, vol. 85, no. 1, pp. 24--45, 2023, doi:10.1093/jrsssb/qkac004.

\bibitem{ref19} K. M. Cohen, S. Park, O. Simeone, and S. Shamai (Shitz), "Calibrating AI models for wireless communications via conformal prediction," arXiv:2212.07775, 2022.

\bibitem{ref20} O. Simeone, S. Park, and M. Zecchin, "Conformal calibration: ensuring the reliability of black-box AI in wireless systems," arXiv:2504.09310, 2025.

\bibitem{ref21} A. N. Angelopoulos, S. Bates, A. Fisch, L. Lei, and T. Schuster, "Conformal risk control," in Proc. International Conference on Learning Representations (ICLR), 2024.

\bibitem{ref22} S. Bates, A. Angelopoulos, L. Lei, J. Malik, and M. I. Jordan, "Distribution-free, risk-controlling prediction sets," Journal of the ACM, vol. 68, no. 6, 2021, doi:10.1145/3478535.

\bibitem{ref23} A. G. Hawkes, "Spectra of some self-exciting and mutually exciting point processes," Biometrika, vol. 58, no. 1, pp. 83--90, 1971.

\bibitem{ref24} P. J. Laub, Y. Lee, P. K. Pollett, and T. Taimre, "Hawkes models and their applications," Annual Review of Statistics and Its Application, vol. 12, pp. 233--258, 2025.

\bibitem{ref25} Y. Ogata, "Statistical models for earthquake occurrences and residual analysis for point processes," Journal of the American Statistical Association, vol. 83, no. 401, pp. 9--27, 1988.

\bibitem{ref26} M. Price-Williams and N. A. Heard, "Nonparametric self-exciting models for computer network traffic," Statistics and Computing, vol. 30, no. 2, pp. 209--220, 2020.

\bibitem{ref27} H. Ismail Fawaz, G. Del Grosso, T. Kerdoncuff, A. Boisbunon, and I. Saffar, "Deep unsupervised domain adaptation for time series classification: a benchmark," arXiv:2312.09857, 2023.

\bibitem{ref28} B. Sun and K. Saenko, "Deep CORAL: correlation alignment for deep domain adaptation," in Computer Vision -- ECCV 2016 Workshops, LNCS vol. 9915, pp. 443--450, 2016.

\bibitem{ref29} Y. Shi, X. Ying, and J. Yang, "Deep unsupervised domain adaptation with time series sensor data: a survey," Sensors, vol. 22, no. 15, art. 5507, 2022.

\bibitem{ref30} D. Wagner, A. Nair, B. J. Franks, et al., "Formally exploring time-series anomaly detection evaluation metrics," arXiv:2510.17562, 2025.

\bibitem{ref31} S. Deb, M. Rathod, R. Balamurugan, S. K. Ghosh, R. K. Singh, and S. Sanyal, "Performance evaluation of conditional handover in 5G systems under fading scenario," arXiv:2403.04379, 2024.

\bibitem{ref32} G. Ke, Q. Meng, T. Finley, T. Wang, W. Chen, W. Ma, Q. Ye, and T.-Y. Liu, "LightGBM: a highly efficient gradient boosting decision tree," in Advances in Neural Information Processing Systems 30, 2017, pp. 3146--3154.

\bibitem{ref33} C. Guo, G. Pleiss, Y. Sun, and K. Q. Weinberger, "On calibration of modern neural networks," in Proc. ICML, 2017, pp. 1321--1330.

\bibitem{ref34} B. Zadrozny and C. Elkan, "Transforming classifier scores into accurate multiclass probability estimates," in Proc. ACM SIGKDD, 2002, pp. 694--699.

\bibitem{ref35} D. R. Roberts et al., "Cross-validation strategies for data with temporal, spatial, hierarchical, or phylogenetic structure," Ecography, vol. 40, no. 8, pp. 913--929, 2017, doi:10.1111/ecog.02881.

\bibitem{ref36} S. Kaufman, S. Rosset, C. Perlich, and O. Stitelman, "Leakage in data mining: formulation, detection, and avoidance," ACM Transactions on Knowledge Discovery from Data, vol. 6, no. 4, art. 15, 2012, doi:10.1145/2382577.2382579.

\bibitem{ref37} D. Raca, D. Leahy, C. J. Sreenan, and J. J. Quinlan, "Beyond throughput, the next generation: a 5G dataset with channel and context metrics," in Proc. ACM Multimedia Systems Conference (MMSys), 2020, pp. 303--308, doi:10.1145/3339825.3394938. Dataset: github.com/uccmisl/5Gdataset (GPL-3.0).

\bibitem{ref38} M. M. S. Shafi, K. M. Istiaque, S. S. Sowad, and M. T. Kawser, "Drive-test-based LTE handover dataset for cellular mobility studies in urban Bangladesh," Mendeley Data, version 2, 2025, doi:10.17632/n2pvmtyn2j.2.

\bibitem{ref39} M. M. S. Shafi, K. M. Istiaque, S. S. Sowad, and M. T. Kawser, "Handover optimization in LTE networks using contextual bandit reinforcement learning and real-world data," Technical report / manuscript, Department of Electrical and Electronic Engineering, Islamic University of Technology, 2025. Available on ResearchGate.

\bibitem{ref40} F. Chen, M. Ghoshal, E. Nan, P. Dinh, I. Khan, Z. J. Kong, Y. C. Hu, and D. Koutsonikolas, "A large-scale study of the potential of multi-carrier access in the 5G era," in Passive and Active Measurement (PAM), LNCS vol. 15567, 2025, pp. 469--484, doi:10.1007/978-3-031-85960-1\_19. Dataset: github.com/NUWiNS/pam2025-multi-carrier-dataset.

\bibitem{ref41} M. Ghoshal, I. Khan, Z. J. Kong, P. Dinh, J. Meng, Y. C. Hu, and D. Koutsonikolas, "Performance of cellular networks on the wheels," in Proc. ACM Internet Measurement Conference (IMC), 2023, doi:10.1145/3618257.3624814.

\bibitem{ref42} D. Hong and S. S. Rappaport, "Traffic model and performance analysis for cellular mobile radio telephone systems with prioritized and nonprioritized handoff procedures," IEEE Transactions on Vehicular Technology, vol. 35, no. 3, pp. 77--92, 1986.

\bibitem{ref43} X. Lin, R. K. Ganti, P. J. Fleming, and J. G. Andrews, "Towards understanding the fundamentals of mobility in cellular networks," IEEE Transactions on Wireless Communications, vol. 12, no. 4, pp. 1686--1698, Apr. 2013, doi:10.1109/TWC.2013.022113.120512.

\bibitem{ref44} S. Sadr and R. S. Adve, "Handoff rate and coverage analysis in multi-tier heterogeneous networks," IEEE Transactions on Wireless Communications, vol. 14, no. 5, pp. 2626--2638, May 2015, doi:10.1109/TWC.2015.2390231.

\bibitem{ref45} M. Kabeer, R. Nordin, M. Behjati, and S. L. Lau, "Sequence-based deep learning for handover optimization in dense urban cellular network," arXiv:2510.02958, 2025. [Preprint; not part of the screened set of Section 2.3.]

\bibitem{ref46} S. V. R. Mandapati, "Meta-continual mobility forecasting for proactive handover prediction," arXiv:2512.11841, 2025. [Preprint; not part of the screened set of Section 2.3.]

\bibitem{ref47} A. Qchohi, J. Moysen Cortes, and M. Zecchin, "Confounding-valid conformal inference for counterfactual KPIs in wireless networks," arXiv:2609.05073, 2026. [Preprint.]

\bibitem{ref48} R. L. Prentice and L. A. Gloeckler, "Regression analysis of grouped survival data with application to breast cancer data," Biometrics, vol. 34, no. 1, pp. 57--67, 1978, doi:10.2307/2529588.

\bibitem{ref49} P. D. Allison, "Discrete-time methods for the analysis of event histories," Sociological Methodology, vol. 13, pp. 61--98, 1982, doi:10.2307/270718.

\bibitem{ref50} G. Tutz and M. Schmid, \textit{Modeling Discrete Time-to-Event Data}, Springer Series in Statistics, Springer, Cham, 2016, doi:10.1007/978-3-319-31552-2.

\bibitem{ref51} R. J. Tibshirani, R. Foygel Barber, E. Cand\`es, and A. Ramdas, "Conformal prediction under covariate shift," in \textit{Advances in Neural Information Processing Systems 32 (NeurIPS)}, 2019, pp. 2530--2540.

\bibitem{ref52} R. Foygel Barber, E. J. Cand\`es, A. Ramdas, and R. J. Tibshirani, "Conformal prediction beyond exchangeability," \textit{The Annals of Statistics}, vol. 51, no. 2, pp. 816--845, 2023, doi:10.1214/23-AOS2276.

\bibitem{ref53} I. Gibbs and E. Cand\`es, "Adaptive conformal inference under distribution shift," in \textit{Advances in Neural Information Processing Systems 34 (NeurIPS)}, 2021, pp. 1660--1672.
\end{thebibliography}
\onehalfspacing
